% !TEX root = ../../技术文档.tex
\section{系统互连与存储}

\subsection{总线结构}

CPU采用32\,bit地址、64\,bit取指返回和32\,bit数据访问接口，经本地存储路由连接Boot ROM、SRAM、NPU、Video及DDR。
应用采用片上程序与数据存储配置，DDR访问通过不可缓存通路接入共享互连。
控制寄存器采用32\,bit AXI4-Lite，数据DMA采用128\,bit AXI4，通道握手与响应遵循AXI协议\cite{arm-axi-ihi0022e}，在DDR封装处适配厂商IP的用户位宽和时钟。
AXI互连按事务所有者串行转发读写请求，完成当前事务后再接纳下一请求。

跨时钟域通路在明确的桥接边界使用异步FIFO或握手机制。运行错误使系统停止接纳新任务，
已接纳的事务和缓冲引用保持归属，重新启动通过协调复位建立统一会话。

\subsection{存储组织}

CPU启动存储包括8\,KiB Boot ROM和32\,KiB SRAM。板载DDR3物理容量为1\,GB，
软件映射其中128\,MiB窗口：32\,MiB划为CPU私有区，96\,MiB作为不可缓存共享区。应用代码、数据及栈均位于SRAM。
系统地址映射按下表组织，地址属性同时用于CPU访存检查、驱动和链接配置。

\begin{reporttable}
\begin{longtable}{@{}p{0.23\linewidth}p{0.26\linewidth}p{0.43\linewidth}@{}}
\caption{系统地址映射表}\\
\toprule
区域 & 基址 & 容量或用途 \\
\midrule
\endfirsthead
\caption[]{系统地址映射表（续）}\\
\toprule
区域 & 基址 & 容量或用途 \\
\midrule
\endhead
\bottomrule
\endfoot
Boot ROM & \texttt{0x0000\,0000} & 8\,KiB，启动代码 \\
片上SRAM & \texttt{0x0001\,0000} & 32\,KiB，应用代码、数据及栈 \\
NPU控制 & \texttt{0x2000\,0000} & 整帧任务、完成记录及状态 \\
Video控制 & \texttt{0x2001\,0000} & 采集、显示提交、释放及状态 \\
CPU私有DDR & \texttt{0x8000\,0000} & 32\,MiB，CPU私有地址区 \\
共享DDR & \texttt{0x8200\,0000} & 96\,MiB，不可缓存 \\
\end{longtable}
\end{reporttable}

共享区中，帧槽A、B、C的基址分别为\texttt{0x8200\,0000}、\texttt{0x8220\,0000}和\texttt{0x8240\,0000}，每槽2\,MiB。
原图位于槽内偏移0，有效数据量为1,440,000\,B，槽内余量按固定地址布局保留。
从\texttt{0x8260\,0000}开始的2\,MiB为头像素材区，区首保存256$\times$256 RGBA8888头像，有效数据量为262,144\,B。

检测输入区从\texttt{0x8280\,0000}开始，三槽基址依次为\texttt{0x8280\,0000}、\texttt{0x8285\,0000}和\texttt{0x828A\,0000}，每槽预留320\,KiB，有效张量为307200\,B。
检测输出区从\texttt{0x8290\,0000}开始，预留2\,MiB。输出槽A、B、C分别从\texttt{0x8290\,0000}、\texttt{0x8294\,0000}和\texttt{0x8298\,0000}开始，每槽256\,KiB。
320$\times$320输入对应2100个预测位置，每位置包含6个16\,bit值，有效头数据为25,200\,B，各头分配跨度合计26,368\,B。
输出头位于槽内偏移0，其余空间按公共布局保留；整帧完成信息通过MMIO寄存器返回，CPU读取输出头形成检测结果。

模型区预留8\,MiB，基址为\texttt{0x82B0\,0000}。NPM3程序占10,240\,B，放在区首；CPU解码表从\texttt{0x82B1\,0000}开始，长度1,572,864\,B。
张量与参数工作区从\texttt{0x8330\,0000}开始，预留64\,MiB；本部署模型实际使用4,876,544\,B。
参数镜像位于工作区偏移4,768,256，长度108,288\,B，保存权重、Bias、重定标及辅助算子参数，前部空间用于中间张量。
共享区余下13\,MiB用于装载、测试和扩展。
跨层张量按生命周期分配，帧槽、输入输出槽及模型工作区分别管理。
CPU提交整帧任务时，以共享窗口基址\texttt{0x8200\,0000}计算程序、工作区及输入输出槽的相对地址；NPU计算图执行器解析模型内的张量引用。
内部存储服务地址相对DDR基址\texttt{0x8000\,0000}，服务适配器在发起总线访问时加该基址一次。

\subsection{访问仲裁}

CPU、NPU和视频DMA共享DDR，互连采用轮询方式选择事务所有者。
仲裁器记录每笔事务的发起主端，并将响应、背压和错误返回到原请求方。

带宽统计分别使用实际采集率和显示扫描率。原图及模型输入的写流量按被接纳的完整帧计数，
全屏显示读流量按原图扫描率及行缓冲的实际读取量计算。

\subsection{DDR带宽预算}

每帧采集写入1,440,000\,B原图和307,200\,B输入张量，共1,747,200\,B。
全屏显示每次扫描读取1,440,000\,B原图；即使重复显示同一帧，该次扫描仍产生读流量。

40\,MHz像素时钟配合1056$\times$628总时序，扫描频率为
\[
f_{\mathrm{scan}}=\frac{40\times10^6}{1056\times628}\approx60.31654\,\mathrm{Hz}
\]
以下带宽预算采用该扫描频率。
设接纳帧率为$f_{\mathrm{in}}$，则视频有效载荷带宽为
\[
B_{\mathrm{video}}=1{,}747{,}200f_{\mathrm{in}}+1{,}440{,}000f_{\mathrm{scan}}\quad\mathrm{B/s}
\]

\begin{reporttable}
\footnotesize
\begin{longtable}{@{}rrrrr@{}}
\caption{两支路采集与全屏显示的DDR有效载荷预算}\\
\toprule
接纳帧率 & 采集写入 & 显示读取 & 读写合计 & 占原生端口峰值 \\
(fps) & (MB/s) & (MB/s) & (MB/s) & (\%) \\
\midrule
\endfirsthead
\caption[]{两支路采集与全屏显示的DDR有效载荷预算（续）}\\
\toprule
接纳帧率 & 采集写入 & 显示读取 & 读写合计 & 占原生端口峰值 \\
(fps) & (MB/s) & (MB/s) & (MB/s) & (\%) \\
\midrule
\endhead
\bottomrule
\endfoot
24 & 41.933 & 86.856 & 128.789 & 3.220 \\
30 & 52.416 & 86.856 & 139.272 & 3.482 \\
60 & 104.832 & 86.856 & 191.688 & 4.792 \\
\end{longtable}
\end{reporttable}

表中MB按$10^6$\,B计，24、30和60\,fps为容量预算场景，实际应用吞吐由完整处理链路决定。
接纳帧率表示进入处理链路的完整帧数，有效输出帧率表示完成检测并形成可用结果的帧数，两者需分别计数。

DDR3控制器采用HMIC\_S IP，用户端口和初始化接口依据厂商指南配置。
DDR3配置为32\,bit、1000\,MT/s，对应500\,MHz外部CK；厂商IP原生用户端口为256\,bit、125\,MHz，
两种口径的理论数据峰值均为4,000\,MB/s。共享AXI4数据通路为128\,bit、125\,MHz，
理想连续传输上限为2,000\,MB/s，以上三种视频载荷分别占其6.439\%、6.964\%和9.584\%。
原生端口的读写共享物理存储带宽，不能把两个方向的峰值相加。
事务串行仲裁、桥接缓冲、写后读回校验、刷新、行切换和读写换向均会增加占用，
因此载荷占比仅用于预算，剩余百分比不是已经测得的可用吞吐余量。
该桥在写响应前读回核验已写数据；仅按采集有效字节等量读回一次计算，三种场景还需增加
41.933、52.416、104.832\,MB/s原生读流量，视频及其写校验的合计下限为
170.721、191.688、296.520\,MB/s，突发对齐与其他开销需继续计入。

NPU读取输入、权重和中间张量并写回输出的流量另计；采集端307,200\,B写入不包含NPU随后读取该输入的流量。
总预算还需加入CPU访存、NPU计算图执行及启动装载流量，不能只以视频载荷判断整机带宽达标。
头像在启动阶段从DDR读入片上纹理缓存，连续显示时由缓存供数；全尺寸素材装载与缓存填充单独计入启动流量，不作为每次扫描的DDR读流量。
应用配置及原始验证记录集中见\path{docs/verification/soc/face_photo_report_20261104/}。

\subsection{共享数据访问约定}

CPU应用代码、栈及检测候选工作区位于片上SRAM，DMA共享张量、帧数据和MMIO使用不可缓存访问。
CPU通过共享DDR读取九头输出和解码表，保持与NPU写回结果一致的数据视图。
写入方完成数据写入并执行所需屏障后发布就绪事件，读取方完成读取后才交还缓冲区。

采集完成等待两条写入分支的成功响应，NPU完成等待整图执行结束、DMA无在途事务及外部服务适配器的发布许可，
显示切换等待安全帧边界，旧槽释放等待所有显示DMA、行缓存预取和FIFO引用排空。
CPU在完成记录保持期间读取检测头并完成后处理，构建显示记录后ACK整帧任务；确认NPU返回IDLE后归还整图工作区，原图槽继续按显示引用管理。
访问属性、总线完成和帧所有权共同决定数据何时可读、可写。

